% ============================================================
%  Part XVI — Everyday German
% ============================================================
\gpart{XVI}{Everyday German}

% ------------------------------------------------------------
\gsec{Greetings and politeness}

\gintro{Formality matters more in German than in English. The safe default with
anyone you do not know is \textit{Sie} and \textit{Guten Tag}; wait to be
offered \textit{du}.}

\begin{dtbl}[\footnotesize]{W{1.05}W{0.95}W{1}}
\rowcolor{Ink} \hd{Deutsch} & \hd{English} & \hd{When} \\ \hdrule
Guten Morgen        & good morning        & until about 11 \\
Guten Tag           & hello               & the neutral default \\
Guten Abend         & good evening        & from about 18:00 \\
Hallo               & hi                  & informal, any time \\
Grüß Gott           & hello               & southern Germany, Austria \\
Servus              & hi, bye             & Bavaria, Austria \\
Auf Wiedersehen     & goodbye             & formal \\
Tschüss             & bye                 & informal \\
Bis bald            & see you soon        & \\
Bis morgen          & see you tomorrow    & \\
Gute Nacht          & good night          & going to bed \\
Wie geht es Ihnen?  & how are you?         & formal \\
Wie geht's?         & how are you?         & informal \\
Danke, gut          & fine, thanks         & \\
Und Ihnen?          & and you?             & formal \\
Bitte               & please               & \\
Danke               & thank you            & \\
Danke schön         & thank you very much  & \\
Bitte schön         & you're welcome       & also \emph{here you are} \\
Gern geschehen      & my pleasure          & \\
Entschuldigung      & excuse me, sorry     & \\
Es tut mir leid     & I'm sorry            & apology \\
Kein Problem        & no problem           & \\
Herzlichen Glückwunsch & congratulations   & \\
Viel Glück          & good luck            & \\
Gute Besserung      & get well soon        & \\
Guten Appetit       & enjoy your meal      & \\
Prost               & cheers               & \\
\end{dtbl}

\begin{gnote}
\textit{Bitte} does the work of four English words: \emph{please},
\emph{you're welcome}, \emph{here you are}, and — as \textit{Wie bitte?} —
\emph{pardon?}
\end{gnote}

% ------------------------------------------------------------
\gsec{Getting by}

\begin{dtbl}[\footnotesize]{W{1.15}W{0.85}W{1}}
\rowcolor{Ink} \hd{Deutsch} & \hd{English} & \hd{Note} \\ \hdrule
Sprechen Sie Englisch?     & do you speak English? & \\
Ich spreche kein Deutsch   & I don't speak German & \\
Ich verstehe nicht         & I don't understand & \\
Können Sie das wiederholen? & could you repeat that? & \\
Langsamer, bitte           & more slowly, please & \\
Wie bitte?                 & pardon? & \\
Was bedeutet das?          & what does that mean? & \\
Wie sagt man \dots\ auf Deutsch? & how do you say \dots\ in German? & \\
Wie heißen Sie?            & what is your name? & formal \\
Ich heiße \dots            & my name is \dots & \\
Ich komme aus \dots        & I'm from \dots & \\
Wo ist \dots?              & where is \dots? & \\
Wie komme ich zum Bahnhof? & how do I get to the station? & \\
Wie weit ist es?           & how far is it? & \\
Was kostet das?            & how much is that? & \\
Ich hätte gern \dots       & I'd like \dots & polite, Konjunktiv II \\
Ich möchte \dots           & I'd like \dots & equally polite \\
Die Rechnung, bitte        & the bill, please & \\
Zahlen, bitte              & the bill, please & shorter, in a bar \\
Haben Sie \dots?           & do you have \dots? & \\
Ich suche \dots            & I'm looking for \dots & \\
Nur schauen, danke         & just looking, thanks & \\
Können Sie mir helfen?     & can you help me? & \\
Ich brauche einen Arzt     & I need a doctor & \\
Mir ist schlecht           & I feel sick & dative \\
Ich habe mich verlaufen    & I'm lost & on foot \\
Wo sind die Toiletten?     & where are the toilets? & \\
Vorsicht!                  & careful! & \\
\end{dtbl}

\begin{gnote}
\textit{Ich hätte gern} and \textit{Ich möchte} are the two standard polite
requests. \textit{Ich will} means \emph{I want} and sounds blunt to the point of
rudeness when ordering.
\end{gnote}
